\section{The Method Under Evaluation}
\label{sec:method}

Part I concluded that the classical reconstruction budget is where an
intervention should go. The intervention evaluated here is
confidence-guided adaptive reconstruction, specified in full in the
companion article~\cite{methodpaper}; this section states only what is
needed to read the results.

A per-window posterior comparison statistic $\Cw$ --- the posterior
probability, under a two-proportion Beta model with Jeffreys priors, that
the most-observed bit pattern's true sampling probability exceeds the
runner-up's --- is computed from counts already measured. It drives three
independently switchable modules, each replacing one of the three fixed
budgets of Sec.~\ref{sec:pipeline}: \textbf{B} retains candidates to a
target probability mass instead of a fixed top-$k$; \textbf{C} scales the
beam width by how many windows needed that expansion; \textbf{D} resamples
windows whose $\Cw$ falls below a threshold, up to a hard cap. With all
three disabled the implementation reduces to the published pipeline, which
Sec.~\ref{sec:baselinecheck} verifies trial by trial.

The quantum circuit, its width, its depth and its sampling procedure are
unchanged. Five arms are compared throughout: Baseline, B, C, D, and Full
(all three enabled).

% =====================================================================
